
All five sub-hypotheses passed validation gates with metrics exceeding targets, demonstrating mechanistic feasibility of the three-component system at proof-of-concept scale.

\subsubsection{5.1 Instrumentation Feasibility (h-e1)}

\textbf{Overhead:} 0.21\% (target: <10\%)  
\textbf{Capture Rate:} 100\% (target: $\geq 95$\%)  
\textbf{Events Captured:} 100 over simulated 6-month window (target: $\geq 100)$

Load-time instrumentation added negligible latency (<2ms per dataset load) across 30 benchmark runs (3 datasets $\times$ 10 runs each). SHA256 hashing was amortized via lazy caching. Async queue writes minimized blocking I/O. Overhead significantly below 10\% tolerance suggests production deployment is viable without performance degradation.

Per-dataset breakdown showed consistent overhead across dataset sizes: CIFAR-10 (max 0.21\%, mean 0.07\%), IMDB (max 0.06\%, mean 0.04\%), WikiText-103 (max 0.08\%, mean 0.05\%). Telemetry transmission success rate was 100\% (30/30 successful writes to SQLite backend).

\subsubsection{5.2 Health Metrics Precision (h-m1)}

\textbf{Precision:} 88.3\% (target: $\geq 60$\%)  
\textbf{Recall:} 100\% (target: $\geq 80$\%)  
\textbf{Weighted Score Threshold:} $\geq 2.5$

Automated health metrics exceeded precision target by 47\%. Weighted scoring (velocity=2.0, emergence=1.5, issue\_ratio=1.0) achieved strong discrimination on 1000-dataset simulation over a 6-month observation window. Zero false negatives (100\% recall) validates that all true deprecation candidates were surfaced by metrics.

Confusion matrix analysis: 181 true positives (deprecated datasets correctly flagged), 24 false positives (active datasets incorrectly flagged), 0 false negatives (no missed deprecations), 795 true negatives (active datasets correctly not flagged). Declining usage velocity (velocity < 1.0) was the strongest predictor. Issue ratio contributed less signal but improved precision for edge cases (maintained datasets with declining usage due to successor availability rather than quality issues).

Threshold sensitivity analysis on velocity thresholds (0.5, 1.0, 1.5) with fixed emergence and issue thresholds showed that velocity=1.0 optimally balanced precision (88.3\%) and recall (100\%).

\subsubsection{5.3 Context Inference Accuracy (h-m2)}

\textbf{Accuracy:} 100\% (target: $\geq 70$\%)  
\textbf{User Override Rate:} 27.95\% (target: <50\%)  
\textbf{Edge Precision:} 100\% (target: $\geq 60$\%)

Pattern-based context inference achieved perfect accuracy on a 500-sample synthetic validation dataset (150 test samples after 70/30 split). The 72.05\% acceptance rate (100\% - 27.95\% override) suggests context-aware recommendations add value over linear version succession.

Breakdown by context type: classification (100\% accuracy via sklearn/xgboost pattern matching), pretraining (100\% via transformers/torchvision pattern matching), robustness (100\% via foolbox/cleverhans pattern matching), unknown fallback (100\%).

Accuracy is expected to regress to 70-90\% on real-world usage patterns because synthetic patterns are cleaner than production data. The override rate provides a safety valve: users can correct misclassified context, preventing incorrect successor recommendations. Citation parsing from dataset cards achieved 100\% edge inference accuracy on patterns such as ''improved version of X'' and ''successor to X.''

\subsubsection{5.4 Migration Planning Completeness (h-m3)}

\textbf{Impact Completeness:} 100\% (target: $\geq 95$\%)  
\textbf{Schema Accuracy:} 100\% (target: $\geq 80$\%)

Transitive closure analysis on a 4-node dependency graph (1 deprecated dataset, 1 successor dataset, 2 dependent entities) identified all affected entities (100\% recall). Field-level schema diff detected all breaking changes (column renames, type changes, missing fields) automatically with 100\% accuracy.

Schema compatibility analysis detected 2 breaking changes: FIELD\_REMOVAL for \texttt{source} field (MAJOR breaking change) and TYPE\_CHANGE for \texttt{label} field from int32 to int64 (MINOR breaking change). Overall compatibility was classified as MAJOR\_BREAKING.

Validated on small-graph scale (4 nodes); 50,000-node scalability is inferred but not tested. NetworkX graph algorithms scale to thousands of nodes with sub-second query times based on O(V+E) complexity analysis. Production validation is required for 60,000-dataset HuggingFace scale. Plan generation latency was <100ms (target: <5s).

\subsubsection{5.5 Full-Stack Telemetry (h-m4)}

\textbf{Capture Rate:} 100\% (95\% CI: 99.05\%-100\%) (target: $\geq 95$\%)  
\textbf{Overhead:} 5\% (target: <10\%)  
\textbf{Event Completeness:} 100\% (target: $\geq 95$\%)

The integrated system (health metrics + context graphs + instrumentation) maintained high capture rate with overhead below tolerance. Async queue design scaled to 100 simulated user sessions (400 total load operations) without blocking. Wilson confidence interval lower bound (99.05\%) exceeds 95\% target, validating statistical reliability.

Overhead was higher than h-e1 baseline (5\% versus 0.21\%) due to the full telemetry stack (metadata capture, dependency graph queries, health metric computation) but remained well below 10\% tolerance. Baseline load latency was 100ms (mean); instrumented load latency was 105ms (mean), yielding 5\% overhead.

Adoption event tracking captured 33 adoption events (expected 20 based on 40\% encounter rate $\times$ 50\% adoption rate $\times$ 100 users). Higher-than-expected adoption events were due to simulation dynamics where users loaded successor datasets multiple times. Capture mechanism tracked all adoption paths correctly.

\subsubsection{5.6 Cross-Hypothesis Patterns}

All five sub-hypotheses (h-e1, h-m1, h-m2, h-m3, h-m4) passed gates, validating mechanistic integration:

\begin{itemize}
\item \textbf{Component 1 (Health Metrics):} Automated detection feasible with 88.3\% precision, 100\% recall.
\item \textbf{Component 2 (Context Graphs):} Task-aware recommendations feasible with 100\% accuracy on synthetic patterns.
\item \textbf{Component 3 (Instrumentation):} Adoption tracking feasible with 0.21-5\% overhead and 100\% capture.
\end{itemize}

Results validate the infrastructure layer (measurement and mechanisms work) but not the efficacy layer (adoption lift untested because Phase 5 baseline comparison was skipped).

\subsubsection{5.7 Surprising Findings}

Instrumentation overhead was significantly lower than expected (0.21\% baseline, 5\% full stack). SHA256 hashing was expected to dominate overhead; lazy caching and background queue writes effectively amortized cost. Perfect scores (100\% recall, 100\% context accuracy) are likely optimistic due to synthetic data. Real-world accuracy is expected to regress to 70-95\% based on pattern cleanliness assumptions. All metrics exceeded targets by 5-40 percentage points, suggesting proof-of-concept validation was successful but may not reflect production-scale edge cases.
